\documentclass[UTF8]{ctexart}

\makeatletter
\def\input@path{{../}{../../}}
\makeatother
\usepackage{researchnotes}

\title{加法定理与容斥原理}
\author{}
\date{}

\begin{document}
\maketitle

\section{定理与证明}
\label{sec:inclusion-exclusion-theorem-proof}

\begin{theorem}[加法定理]
\label{thm:inclusion-exclusion}
设 $A_1,\ldots,A_n$ 是同一概率空间中的事件，$n\ge2$，则
\begin{equation}
  P\left(\bigcup_{i=1}^{n}A_i\right)
  =\sum_{r=1}^{n}(-1)^{r-1}
    \sum_{1\le i_1<\cdots<i_r\le n}
    P\left(A_{i_1}\cap\cdots\cap A_{i_r}\right).
  \label{eq:inclusion-exclusion-general}
\end{equation}
此式称为概率的\keyterm{加法定理}（addition theorem），也称\keyterm{容斥原理}（inclusion--exclusion principle）。内层求和遍历所有由 $r$ 个不同下标确定的交集，每组指标只取一次。因此，公式依次为单事件概率之和，减去两两交集概率之和，加上三者交集概率之和，直至符号为 $(-1)^{n-1}$ 的全体事件交集项；不要求事件相互独立。
\end{theorem}

\begin{proof}
采用\keyterm{数学归纳法}（mathematical induction）。当 $n=2$ 时，$A_1\cup A_2$ 是 $A_1$ 与 $A_2\setminus A_1$ 的不交并，而 $A_2$ 是 $A_2\setminus A_1$ 与 $A_1\cap A_2$ 的不交并。由概率的有限可加性，
\begin{equation}
  \begin{aligned}
    P(A_1\cup A_2)
    &=P(A_1)+P(A_2\setminus A_1)\\
    &=P(A_1)+P(A_2)-P(A_1\cap A_2).
  \end{aligned}
  \label{eq:inclusion-exclusion-two}
\end{equation}
这正是两个事件的公式，其中 $E\setminus F$ 表示属于 $E$ 而不属于 $F$ 的部分。

设 $k\ge3$，并假设式~\eqref{eq:inclusion-exclusion-general} 对任意 $k-1$ 个事件成立。令 $U=\bigcup_{i=1}^{k-1}A_i$，$B_i=A_i\cap A_k$（$1\le i\le k-1$）。由集合的分配律，$U\cap A_k=\bigcup_{i=1}^{k-1}B_i$，再用二事件公式~\eqref{eq:inclusion-exclusion-two}，得
\begin{equation}
  P\left(\bigcup_{i=1}^{k}A_i\right)
  =P(U)+P(A_k)-P\left(\bigcup_{i=1}^{k-1}B_i\right).
  \label{eq:inclusion-exclusion-induction}
\end{equation}

分别对 $A_1,\ldots,A_{k-1}$ 和 $B_1,\ldots,B_{k-1}$ 应用归纳假设。式~\eqref{eq:inclusion-exclusion-induction} 中的 $P(U)$ 给出所有不含 $A_k$ 的交集项，$P(A_k)$ 给出只含 $A_k$ 的一阶项。最后一项的展开中，每个 $r$ 重交集满足
\[
  B_{i_1}\cap\cdots\cap B_{i_r}
  =A_{i_1}\cap\cdots\cap A_{i_r}\cap A_k,
  \qquad 1\le r\le k-1.
\]
它对应一个含 $A_k$ 的 $r+1$ 重交集；归纳公式中的系数为 $(-1)^{r-1}$，再乘以外面的负号，变为 $(-1)^r$，恰好是 $r+1$ 重交集应有的符号。每组指标唯一确定一个交集项，因此所有交集项都不重不漏地出现，合并后即为式~\eqref{eq:inclusion-exclusion-general} 在 $n=k$ 时的形式。归纳完成。
\end{proof}

\end{document}
